% EXAMPLE_REPORT.tex — extended example for CA22
\documentclass[11pt]{article}
\usepackage[utf8]{inputenc}
\usepackage{geometry}
\geometry{margin=1in}
\usepackage{amsmath,amssymb}
\usepackage{graphicx}
\usepackage{booktabs}
\usepackage{caption}
\usepackage[round]{natbib}
\usepackage{hyperref}
\hypersetup{colorlinks=true,linkcolor=blue,citecolor=blue,urlcolor=blue}

\title{Constrained Policy Optimization on a Synthetic Control Problem}
\author{Example Author \\\texttt{author@example.com}}
\date{\today}

\begin{document}
\maketitle

\begin{abstract}
We study a simple policy-gradient agent on a synthetic episodic control task with an additional constraint cost (a penalty-like "safety" signal). Using a Lagrangian formulation, we seek to maintain expected cost below a threshold while maximizing reward. Our results on a small synthetic dataset demonstrate effective control of constraint expectations with modest impact on reward.
\end{abstract}

\section{Introduction}
Constrained reinforcement learning is important when deploying agents in domains with safety or resource constraints. We experiment with a compact synthetic environment designed to be reproducible and interpretable, allowing us to focus on algorithmic behavior rather than environment engineering.

\section{Background and Related Work}
Classic references on RL include \citet{Sutton2018}, and constrained or safe-RL has recent treatments in the literature (e.g., Lagrangian-based methods; \citet{Achiam2017}). We implement a basic Lagrange multiplier approach as a pedagogical baseline.

\section{Methods}
\subsection{Model Architectures}
We use a small MLP policy and a value baseline (two hidden layers; ReLU activations). See `src/model.py` for code.

\subsection{Objective}
The objective is the standard policy gradient (REINFORCE/actor-critic style) combined with a penalty term of the form:
\begin{equation}
\mathcal{L} = -\mathbb{E}[\log \pi(a|s) \hat{A}] + \mu ( \mathbb{E}[c(s,a)] - c_0 )
\end{equation}
where $\mu$ is the Lagrange multiplier updated by projected gradient ascent (see `src/utils.py`).

\subsection{Dataset}
We use the provided `SyntheticDataset` (short episodes) and perform multiple sampled batches per training step.

\section{Experimental Setup}
We ran small-scale experiments using `configs/debug.yaml` for quick iteration. Key hyperparameters are summarized in Table~\ref{tab:hyper}.

\begin{table}[ht]
\centering
\begin{tabular}{lcc}
\toprule
Hyperparameter & Value \\\midrule
Seed & 0 \\
Hidden size & 16 \\
Batch size & 8 \\
Learning rate & 1e-3 \\
Lagrange lr & 1e-2 \\
Constraint threshold & 0.5 \\
\bottomrule
\end{tabular}
\caption{Main hyperparameters used in the debug experiments.}
\label{tab:hyper}
\end{table}

\section{Results}
Figure~\ref{fig:reward} shows a sample learning curve for reward and constraint mean; Table~\ref{tab:results} reports final means over 5 seeds.

\begin{figure}[ht]
  \centering
  \fbox{\parbox[c][4cm][c]{0.9\linewidth}{\centering Placeholder for learning curve (save as outputs/figures/learning_curve.png)}}
  \caption{Learning curve placeholder (reward and constraint mean).}
  \label{fig:reward}
\end{figure}

\begin{table}[ht]
\centering
\begin{tabular}{lcc}
\toprule
Method & Mean Reward & Mean Constraint \\\midrule
Baseline (no penalty) & 1.35 $\pm$ 0.12 & 0.82 $\pm$ 0.04 \\
Lagrangian (this work) & 1.20 $\pm$ 0.10 & 0.48 $\pm$ 0.03 \\
\bottomrule
\end{tabular}
\caption{Summary statistics across seeds (mean $\pm$ standard error).}
\label{tab:results}
\end{table}

\section{Discussion}
The Lagrangian approach successfully reduced expected constraint to near the target with a small tradeoff in reward. The synthetic setting is useful for reproducibility; future work should test larger-scale environments and alternative constraint optimization strategies.

\section*{Reproducibility and Commands}
To compile this report and reproduce the figures locally (after running the demo notebook to generate figures):
\begin{verbatim}
# build example report
./reports/build_report.sh reports/EXAMPLE_REPORT.tex
\end{verbatim}
To run the demo notebook, open `notebooks/demo.ipynb` and run the cells using `configs/debug.yaml`.

\section*{Acknowledgements}
Thanks to course staff and peers for feedback.

\bibliographystyle{plainnat}
\bibliography{refs}

\end{document}
